\section{Results}

\subsection{Frustration-guided nonlocal moves yield target-dependent, auditable policies}

We first tested whether nonlocal navigation can be made \textbf{self-adaptive} without model training, using an error-geometry proxy $\mathrm{Imp}(p;D)$ to select candidate nonlocal moves. In biophysical terms, $\mathrm{Imp}$ plays the role of a landscape \emph{roughness/frustration} proxy derived from residual geometry: it is not a physical energy, but a protocol-fixed scalar field that shapes exploration. On N=16 contact-map proxy tasks across several targets, the $\mathrm{Imp}$-guided selector achieves competitive best Contact\_F1 while exhibiting target-dependent move usage. In particular, the 2CI2 proxy concentrates selection on \textbf{Expander\_Hot} macro-patches, consistent with the interpretation that its residual is dominated by localized repair needs, whereas 1UBQ uses predominantly local moves and global expander patches.

These behaviours are emitted directly in the audit digest (counts\_total / counts\_early / counts\_late), allowing mechanistic A/B tests on the operator policy itself.

\begin{figure}[t]
  \centering
  \maybeincludegraphics[width=\linewidth]{assets/figures/fig2}
  \caption{Frustration-guided selector: aggregated move-usage counts by target (N=16 proxy). 2CI2 shows strong preference for Expander\_Hot macro moves.}
  \label{fig:2}
\end{figure}

\begin{table}[t]
  \centering
\caption{Frustration-guided benchmarks (N=16 proxy): representative top-performing variants.}
  \label{tab:1}
  \maybeincludetable{assets/tables/table1.tex}
\end{table}

Across this proxy suite, the best mean Contact\_F1 reaches 0.962 on 2CI2\_16 and 5PTI\_16 (Imp-Adaptive-v2) and 0.916 on 1UBQ\_16 (Mix-Expander-Hot), while exhibiting different macro-move usage patterns (Fig.~\ref{fig:2}; Table~\ref{tab:1}). Operationally, these nonlocal moves act like \emph{basin-hopping}: structured, auditable rewrites that help traverse barriers induced by low-resolution objectives.

\subsection{Low-resolution objectives induce a measurable hallucination gap}

To expose a failure mode that is invisible to low-resolution objectives, we intentionally optimized Contact\_F1 alone on real proteins truncated to N=24 residues (1UBQ, 2CI2, 5PTI). Contact\_F1 acts as the optimized readout, whereas distogram RMSE and TM-score are treated as hidden audits. This creates a ‘hallucination gap’: configurations achieving high Contact\_F1 while remaining geometrically inconsistent under higher-resolution residuals. In this regime, purely navigation-driven search can converge to high-F1 states that are off-manifold when inspected by distogram RMSE or TM-score.

\begin{figure}[t]
  \centering
  \maybeincludegraphics[width=\linewidth]{assets/figures/fig3}
  \caption{Hallucination gap under Contact-only optimization (N=24). Points with high Contact\_F1 can exhibit low TM-score, motivating higher-resolution semantic and physical gates.}
  \label{fig:3}
\end{figure}

\subsection{From hard phason rejection to projector-style repair}

Early experiments showed that naive hard phason rejection can freeze search when (i) the 6D alphabet is too coarse to represent the target geometry, or (ii) the slice offset $w_0$ is misaligned, causing physically plausible states to be incorrectly labelled ‘non-crystalline’. To address this, we replaced the filter-style gate \texttt{if ph\_max>theta: reject()} with a projector-style operator \texttt{if ph\_max>theta: project\_to\_window()}.

The projector performs a small inner-loop optimization of $w_0$ (a perpendicular-space translation) to reduce phason violation and accept the repaired candidate if it becomes feasible. This converts phason from a passive constraint into an active dynamical relaxation (a phason-flip analogue).

\subsection{Multi-scale symbolic certificates disambiguate ``theory vs protocol'' at 1000-PDB scale}

At 1000-PDB scale, scalar phason proxies such as $\mathrm{ph}_{\mathrm{rms}}$ or $\mathrm{ph}_{\max}$ can be \emph{too low-resolution} to distinguish a genuine theoretical failure from an \emph{observation-protocol} failure (e.g. misaligned slicing axis or offset definition). We therefore introduce a \textbf{multi-scale symbolic certificate} family: instead of a single scalar, we compute distributional, multi-$m$ signatures from stabilized type streams.

\textbf{Protocol.} For each protein, we derive a binary readout $\rho(t)\in\{0,1\}$, then compute length-$m$ sliding windows and stabilize them via $\mathrm{Fold}_m$ (mapping arbitrary bit windows to golden-mean legal types with no adjacent 11), producing a stable-type stream. We report two complementary readouts:
(i) \textbf{$\rho_A$ (geometry certifier)} from perpendicular-space $y^\perp$ using a \textbf{shared-PCA axis} computed once from the real chain and reused across nulls, with robust median thresholding; and
(ii) \textbf{$\rho_B$ (logic certifier)} from the 6D step stream $\Delta n$ using a parity/majority-sign binarization. To avoid false positives from trivial marginals, we compare against strong correlation-destroying nulls (bond-direction shuffle and block-shuffle) in addition to random chains.

\textbf{Full1k results.} On a reproducible set of 1000 PDB chains (length-filtered to $60\le N\le 350$), these multi-scale certificates produce stable, large effect sizes across $m\in\{6,8,10\}$. The shared-PCA $\rho_A$ readout (perpendicular-space geometric order) yields a strong separation (lower type entropy and lower entropy-rate proxy) relative to shuffle and block-shuffle nulls, while the parity $\rho_B$ readout (step-parity/dynamical order) provides complementary evidence that is particularly strong on the entropy-rate proxy (Table~\ref{tab:auric_full1k}). Notably, cross-readout agreement remains close to zero ($\mathrm{corr}(z_A,z_B)\approx 0$), indicating that the two certificates capture distinct aspects of order rather than redundant views of the same signal.

\textbf{Hard-negative decoy control.} To address the reviewer concern that random or shuffle nulls are ``too easy'' (i.e., trivially non-protein-like), we evaluated the geometry-based certificate on classic hard-negative benchmarks from Decoys 'R' Us: \texttt{4state\_reduced} and \texttt{LMDS}, which contain near-native but incorrect folds.

\paragraph{Robustness across decoy families.}
Using the same hybrid protocol and a protocol-fixed \emph{native-derived} axis per target (the axis $u$ is computed from the native $y^\perp$ once and then reused for all decoys), $\rho_A$ type-entropy typically ranks the native structure among the lower-entropy tail of its decoy set, with several targets showing near-complete separation (e.g., 1R69 and 2CRO with Cliff's $\delta \approx -0.98$).
On \texttt{4state\_reduced} (7 targets), the Shared-PCA $\rho_A$ certificate places the native within the lowest 10\% of entropy for 2/7 targets; switching to the stricter \textbf{Best-Native-Axis} protocol improves this separation (3/7 targets in the lowest 10\%).
Beyond \texttt{4state\_reduced}, on the \texttt{LMDS} decoy set (9 targets), the Shared-PCA $\rho_A$ certificate successfully ranks the native structure within the lowest 10\% of entropy for 4/9 targets. Crucially, switching to the stricter \textbf{Best-Native-Axis} protocol \textit{improves} this separation (5/9 targets in the lowest 10\%), demonstrating that the signal arises from intrinsic native order, not protocol artifacts.

\paragraph{Geometric mimicry in hard cases.}
While separation is strong for most targets, specific targets such as \texttt{1CTF} remain challenging across both benchmarks: under Shared-PCA, \texttt{1CTF} exhibits weak or even reversed separation ($\delta>0$), and although Best-Native-Axis restores the expected ordering ($\delta<0$), the native percentile remains relatively high. We characterize \texttt{1CTF} as a case of \textit{geometric mimicry}, where high-quality decoys can preserve coarse-grained certificate signatures despite topological errors. This highlights the role of the geometry-based certificate as a necessary \textit{feasibility gate} rather than a sufficient atomic-level discriminator, motivating its combination with local energy refinement in the search phase.

Full decoy reports and per-target histograms are provided in the accompanying evidence repository.

\paragraph{External server baseline with strict homology exclusions.}
To complement certificate-based controls with an interpretable cross-method agreement check, we prepared a small benchmark comparing QUARK and I-TASSER on the same six short targets used above, and added within-method ablations that exclude close homologous fragments/templates. This provides an external sanity check on whether the predicted global topology is robust to homology leakage and whether two independent pipelines converge to a similar fold class (TM-score $\ge 0.5$ indicates shared topology). Inputs, job tracking, and the TM-align summary template are provided in \texttt{docs/runs/quark\_itasser\_homology\_ablation/}; per-target outcomes will be summarized in Supplementary Table~\ref{tab:quark_itasser_homology_ablation}.

\paragraph{Physical consistency audit (post-hoc).}
As an additional low-cost sanity check that does not modify search, we computed coarse-grained C$\alpha$ physical proxies (bond-length MSE and a non-neighbour steric-clash proxy) and compared them against the geometry-based certificate. On the \texttt{4state\_reduced} hard-negative decoy benchmark, these signals are weakly correlated (Pearson $r\approx 0$), consistent with an \textbf{orthogonal-defense} interpretation: many decoys satisfy local physical constraints while violating global geometric manifold constraints (high geometric type entropy), and a smaller subset exhibits the converse (low geometric type entropy but poor local physics). Figure~S\ref{fig:physics_audit} summarizes this ``dual-filter funnel'' behavior.

\textbf{AlphaFold transparency context (non-accuracy).} To contextualize PhasonFold's auditability claim against widely used predictors, we also include a qualitative transparency comparison against \texttt{AlphaFold} DB (AFDB), contrasting a static AFDB output + pLDDT profile with PhasonFold's replayable audit trajectory (Fig.~\ref{fig:af_vs_omega_transparency}).

\textbf{Blind-sprint feasibility.} Finally, we tested whether the certificates can be used as \emph{active scoring signals} (not just post-hoc discriminators). On short targets (N<100), we ran blind-sprint prototypes that do not use oracle direction guidance and treat TM-score as evaluation-only, optimizing native-derived contact/distogram objectives augmented by certificate penalties. We report two variants:
\begin{itemize}[nosep,leftmargin=*]
  \item \emph{Contact beam-search + certificates}: On 1CRN (N=46), a contact-only beam-search with certificate scoring reaches TM$\approx$0.74 (see the evidence repository under \texttt{docs/runs/}).
  \item \emph{Sparse-distogram MDS + certificate reranking}: As an upper-bound feasibility test using native-derived pairwise distances (300 sampled pairs per target), we ran a suite of 10 short chains (N=46--76, 5 seeds each). Best-of-seeds results: 9/10 targets achieve TM$>$0.98, with 5 reaching TM=1.0 (1CRN, 1R69, 2CRO, 4PTI, 5PTI); one outlier (1SN3) yields TM=0.53 due to missing structure data (see the evidence repository under \texttt{docs/runs/}).
\end{itemize}
These prototypes are not yet the full PhasonFold move stack, and the distogram MDS variant uses native distance information that would not be available in a truly blind setting. Nevertheless, they establish feasibility that certificate-style audit signals can steer search toward higher-quality structures when combined with informative distance priors.

\paragraph{Blind Sprint ablation: $\rho_B$ rescues hard cases, $\rho_A$ can be risky.}
To turn Blind Sprint from a demonstration into a quantitative experiment, we ran an ablation suite on the 7 classic \texttt{4state\_reduced} short targets (N<100) with 5 random seeds each. We compare a baseline contact-driven beam search against three certificate variants: $\rho_A$ only (geometry), $\rho_B$ only (logic/dynamics; velocity readout), and $\rho_A+\rho_B$ (Fig.~\ref{fig:blind_sprint_ablation}).

\textbf{Certificates rescue search failure cases.}
The baseline policy achieves a high best-of-seeds success rate (6/7 targets with TM$>$0.6) but fails on \texttt{1SN3} (best TM=0.31). Enabling the logic/dynamics certificate $\rho_B$ increases robustness to 7/7 (TM$>$0.6 on all targets), boosting \texttt{1SN3} to best TM=0.84. This shows that $\rho_B$ can act as an actionable feasibility signal during blind search rather than a post-hoc diagnostic.

\textbf{Orthogonality and protocol sensitivity of $\rho_A$ vs $\rho_B$.}
On \texttt{3ICB} we observe a sharp divergence: while $\rho_B$ maintains or improves performance (best TM$\approx$0.79), $\rho_A$ alone can guide search into a failure basin (best TM$\approx$0.22), consistent with an over-constrained geometric manifold early in the search. Together, these ablations support our recommended protocol: use $\rho_B$ for search guidance (high acceptance rate) and reserve $\rho_A$ for geometric auditing and decoy discrimination, where it is strongest.

\begin{figure}[t]
  \centering
  \maybeincludegraphics[width=\linewidth]{assets/figures/fig4_search_ablation}
  \caption{\textbf{Search Ablation (Blind Sprint on 4state targets, 5 seeds each).}
  \textbf{(A)} Best-of-seeds TM-score per target for baseline (grey), $\rho_A$ (blue), $\rho_B$-vel (orange), and $\rho_A+\rho_B$ (green). \texttt{1SN3} is rescued by $\rho_B$ (``Rescue''), while \texttt{3ICB} reveals a ``Geometric Trap'' where $\rho_A$ alone fails.
  \textbf{(B)} Success rate (fraction of targets with best-of-seeds TM$>$0.6): baseline 6/7 vs $\rho_B$-vel 7/7.
  \textbf{(C)} TM ECDF over all seeds from the same suite, showing a distributional shift with certificate terms.
  }
  \label{fig:blind_sprint_ablation}
\end{figure}

\begin{table}[t]
  \centering
  \caption{Full1k hybrid protocol (triple232; $m=8$): effect sizes vs strong nulls. Cliff's $\delta<0$ indicates the real chain is more ordered (lower metric) than the null. Percentiles are native percentiles among shuffle replicates.}
  \label{tab:auric_full1k}
  {%
    \setlength{\tabcolsep}{4pt}%
    \renewcommand{\arraystretch}{1.08}%
    \small%
    \resizebox{\linewidth}{!}{%
      \begin{tabular}{lllllll}
        \toprule
        Readout & Metric & mean $\delta$(real,shuffle) & mean $\delta$(real,blk8) & median pct(real,shuffle) \\
        \midrule
        $\rho_A$ (PCA) & $H(\mathrm{type})$ & -0.274 & -0.242 & 0.300 \\
        $\rho_B$ (parity) & $H(\mathrm{type})$ & -0.366 & -0.119 & 0.200 \\
        $\rho_A$ (PCA) & $\widehat{h}_{\mathrm{SMB}}$ & -0.298 & -0.278 & 0.300 \\
        $\rho_B$ (parity) & $\widehat{h}_{\mathrm{SMB}}$ & -0.494 & -0.326 & 0.100 \\
        \bottomrule
      \end{tabular}%
    }%
  }%
\end{table}

\begin{figure}[t]
  \centering
  \begin{subfigure}[t]{0.49\linewidth}
    \centering
    \maybeincludegraphics[width=\linewidth]{../docs/runs/full1k_hybrid_auric/plots/auric_entropy_multiscale_triple232_rhoA_n1000_seed0_triple232_hybrid_pcaA_parityB_shuffle10_blk8x5_r1}
    \caption{$\rho_A$ (shared-PCA): multi-$m$ trends.}
  \end{subfigure}
  \hfill
  \begin{subfigure}[t]{0.49\linewidth}
    \centering
    \maybeincludegraphics[width=\linewidth]{../docs/runs/full1k_hybrid_auric/plots/auric_entropy_multiscale_triple232_rhoB_n1000_seed0_triple232_hybrid_pcaA_parityB_shuffle10_blk8x5_r1}
    \caption{$\rho_B$ (parity): multi-$m$ trends.}
  \end{subfigure}
  \caption{Multi-scale symbolic certificate on full1k under the hybrid protocol. Multi-scale curves summarize median $H(\mathrm{type})$ across $m\in\{6,8,10\}$ for real and null families.}
  \label{fig:auric_multiscale_full1k}
\end{figure}

\subsection{Projector repair enables feasible hard gates on a long-chain $\beta$-barrel (1EMA/GFP)}

On the long-chain GFP target (1EMA, N=221), the projector substantially reduces phason ratio while maintaining or improving structural metrics, whereas filter-style hard rejection incurs a large reject count and can lower TM-score. These results support the interpretation that feasibility at tight gates is limited by alignment ($w_0$) and expressivity, not by the usefulness of phason as a geometric regularizer.

\begin{figure}[t]
  \centering
  \maybeincludegraphics[width=\linewidth]{assets/figures/fig4}
  \caption{Phason ratio vs TM-score for 1EMA (GFP) comparing baseline (no gate), filter-style hard reject, and projector repair. Projector reduces phason violations without freezing exploration.}
  \label{fig:4}
\end{figure}

\subsection{Physical audit: Omega backbones occupy Rosetta energy basins}

Omega v1 outputs in this repository are primarily \textbf{C$\alpha$-only} traces. To assess whether these geometrically generated backbones are \textbf{physically realizable}---i.e., whether they lie within attractive basins of a standard all-atom force field---we performed a constrained relaxation audit using \textbf{Rosetta FastRelax} as an external ``physical auditor''.

Across 6 homology-ablation targets (\texttt{1R69}, \texttt{2CRO}, \texttt{4PTI}, \texttt{1CTF}, \texttt{1DTK}, \texttt{1SHF-A}), the best-scoring relaxed models (20 trajectories per target; see Methods) exhibited \textbf{low structural drift} relative to the original Omega C$\alpha$ traces, confirming that Omega-generated backbones are compatible with Rosetta's physical energy landscape. We additionally observed target-specific accessibility signatures (e.g., \texttt{1CTF} showed a small median--best gap, whereas \texttt{4PTI} exhibited a deep but narrow basin). Full per-target audit results are provided in Supplementary and in \texttt{docs/runs/quark\_itasser\_homology\_ablation/rosetta\_relax\_summary.csv}.

\textbf{Multi-target three-way audit (1R69, 2CRO).} To contextualize this ``physics audit'' against specialist ab initio predictors, we applied the same constrained FastRelax audit to committable server models from \textbf{QUARK} and \textbf{I-TASSER} (model1; $nstruct=20$). TM-scores in Fig.~\ref{fig:5}A use the repository's \textbf{Kabsch single-pass TM-score} (see Supplementary; \texttt{docs/runs/quark\_itasser\_homology\_ablation/server\_summaries.csv}) as a consistent topological readout across targets. While QUARK/I-TASSER achieve higher TM on these server models, the physical audit reveals a complementary advantage: \textbf{Omega exhibits a lower drift distribution under Rosetta relaxation across both targets} (Fig.~\ref{fig:5}B) and reaches a lower Rosetta energy floor under the same restraint strength.

\begin{figure}[t]
  \centering
  \maybeincludegraphics[width=\linewidth]{assets/figures/fig5_multi_target}
  \caption{\textbf{Multi-target three-way physical audit (Omega vs QUARK vs I-TASSER).}
  \textbf{(A)} Kabsch single-pass TM-score to native for model1 on 1R69 and 2CRO.
  \textbf{(B)} Drift distributions under coordinate-constrained FastRelax ($nstruct=20$): C$\alpha$ RMSD to each method's input trace; Omega shows lower drift across both targets.
  \textbf{(C)} Qualitative C$\alpha$ overlay vs native for 1R69.}
  \label{fig:5}
\end{figure}

\subsection{A feasible-tight hard gate $\theta^\star$ makes projector repair necessary across folds}

To test whether phason constraints are merely optional regularizers or enforceable dynamical conditions, we introduced $\theta^\star$ (theta\_star): a per-target feasible-tight threshold derived from a counterfactual feasibility floor. Operationally, $\theta^\star$ is set just above the smallest ph\_max achievable by projector repair, making it tight but satisfiable. In this regime, filter-style hard rejection fails to ever enter the window, while projector repair reaches feasibility and stays within the gate.

\begin{table}[t]
  \centering
  \caption{$\theta^\star$ hard-gate feasibility: filter vs projector repair.}
  \label{tab:2}
  \maybeincludetable{assets/tables/table2.tex}
\end{table}

Concretely, under $\theta^\star$ the filter gate never reaches feasibility on either 1AKE or 1TIM (pass\_final=False), while the projector enters the window after \(\approx 70\) steps (t\_first\_pass 70.5–73.5) with a small number of repair events (phason\_relax\_n 7–12) (Table~\ref{tab:2}).

\begin{figure}[t]
  \centering
  \maybeincludegraphics[width=\linewidth]{assets/figures/fig6}
  \caption{1AKE ph\_max(t) under $\theta^\star$ with horizontal $\theta^\star$ line: filter fails to cross; projector enters window.}
  \label{fig:6}
\end{figure}

\begin{figure}[t]
  \centering
  \maybeincludegraphics[width=\linewidth]{assets/figures/fig7}
  \caption{1TIM ph\_max(t) under $\theta^\star$ with horizontal $\theta^\star$ line: filter fails; projector crosses and remains feasible.}
  \label{fig:7}
\end{figure}

\subsection{A 6D alphabet hierarchy reveals a >0.9 TM-score ceiling (Triple232)}

Hard geometric constraints alone cannot yield high structural accuracy if the discrete move set cannot represent the target. We therefore evaluated a hierarchy of 6D step alphabets—axis12 (±$e_i$), pair72 (±$e_i$±$e_j$) and triple232 (±$e_i$±$e_j$±$e_k$)—using an oracle direction codec that quantizes native C$\alpha$ bond directions. This isolates representational expressivity from global search.

Triple232 raises the oracle TM-score ceiling to \textasciitilde{}0.94 for 1AKE and \textasciitilde{}0.92 for 2O5P tail250. For 1TIM, the full-length ceiling is \textasciitilde{}0.81, but local 100-residue windows exceed 0.92, suggesting that long-chain accuracy requires multi-scale composition rather than a single global slice.

\begin{figure}[t]
  \centering
  \maybeincludegraphics[width=\linewidth]{assets/figures/fig8}
  \caption{Oracle expressivity of axis12, pair72 and triple232 alphabets. Triple232 enables TM-score >0.9 on 1AKE and 2O5P tail250.}
  \label{fig:8}
\end{figure}

\subsection{Sprint-to-0.9: hotspot Triple232 + piecewise projector reaches TM-score >0.9 on 1TIM and a membrane-protein segment}

Finally, we combined the two key enablers—expressivity (Triple232) and dynamical feasibility ($\theta^\star$ + projector repair)—into a targeted A/B runner.

Configuration A (Baseline) uses pair72 directions with a single global projector under $\theta^\star$. Configuration B (Challenger) uses hotspot-only Triple232 re-encoding (to avoid unnecessary branching) together with a piecewise projector that allows $w_0$ to drift across segments, reducing long-range accumulation of phason strain.

This Sprint-to-0.9 runner is an \emph{oracle-guided compilation} stress test: the ``wish'' is perfect (native alignment supplies hotspot locations), and the goal is to verify that the operator stack (patch + projector) can \emph{realize} a 0.9+ TM-score endpoint when representational and feasibility bottlenecks are addressed (see \texttt{artifacts/reports/sprint0p9\_AB\_runner\_report.md}).

On 1TIM (N=247) and 2O5P tail250 (N=250), the Challenger reaches TM-score >0.9 while maintaining $\theta^\star$ feasibility, and also maintains a hydrophobic-belt packing signal on the membrane segment (Fig.~\ref{fig:12}).

\begin{table}[t]
  \centering
  \caption{Sprint-to-0.9 A/B outcomes (baseline vs challenger).}
  \label{tab:3}
  \maybeincludetable{assets/tables/table3.tex}
\end{table}

In particular, on 1TIM the baseline stalls at TM=0.676, whereas the Challenger reaches TM=0.928--0.940 across three seeds (mean 0.936). On 2O5P tail250, the baseline achieves TM=0.703, whereas the Challenger reaches TM=0.917--0.923 (mean 0.921) across three seeds (Table~\ref{tab:3}).

\begin{figure}[t]
  \centering
  \maybeincludegraphics[width=\linewidth]{assets/figures/fig9}
  \caption{TM-score distributions for Sprint-to-0.9 A/B on 1TIM and 2O5P tail250.}
  \label{fig:9}
\end{figure}

\begin{figure}[t]
  \centering
  \maybeincludegraphics[width=\linewidth]{assets/figures/fig10}
  \caption{Phason–TM scatter with $\theta^\star$ horizontal lines. Challenger points remain below $\theta^\star$ while improving TM-score.}
  \label{fig:10}
\end{figure}

\begin{figure}[t]
  \centering
  \maybeincludegraphics[width=\linewidth]{assets/figures/fig11}
  \caption{1TIM phason-strain along the chain. Piecewise projector suppresses long-range accumulation relative to a global slice.}
  \label{fig:11}
\end{figure}

\begin{figure}[t]
  \centering
  \maybeincludegraphics[width=\linewidth]{assets/figures/fig12}
  \caption{2O5P tail250 hydrophobic belt pattern for native vs A/B reconstructions.}
  \label{fig:12}
\end{figure}

\subsection{Blind Sprint: removing oracle TM-score and native membrane labels}

To test whether PhasonFold retains an advantage without directly optimizing oracle TM-score and native-derived membrane annotations, we ran a Blind Sprint protocol that removes these oracle terms from the acceptance objective and replaces them with proxy guides (semantic residuals such as distogram RMSE and a sequence-only hydropathy belt prior; see \texttt{artifacts/reports/blind\_sprint\_AB\_report\_v1.md}).

\begin{table}[t]
  \centering
  \caption{Blind Sprint A/B summary (evaluation TM-score reported; not used for acceptance).}
  \label{tab:4}
  \maybeincludetable{assets/tables/table4.tex}
\end{table}

\begin{figure}[t]
  \centering
  \maybeincludegraphics[width=0.9\linewidth]{../figures/blind\_sprint\_tmscore\_distribution}
  \caption{Blind Sprint TM-score distribution (evaluation-only). Challenger (B) shifts the TM-score distribution upwards on \texttt{1TIM} and \texttt{2O5P} tail250.}
  \label{fig:13}
\end{figure}
